% Part:sets-functions-relations
% Chapter: size-of-sets
% Section: enumerations-alt

\documentclass[../../../include/open-logic-section]{subfiles}

\begin{document}

\olfileid[tr]{sfr}{siz}{enm-alt}

\olsection{Numaralandırmalar ve \usetoken{S}{enumerable} Kümeler}

\begin{editorial}
  Bu kesim, küme teorisinde yapıldığı gibi numaralandırmaları $\Nat$'ın
  (başlangıç parçalarıyla) bire bir örten eşlemeler olarak tanımlar. Bu
  nedenle \olref[enm]{sec}'teki tanımlarla biraz çatışır ve oradaki bütün
  örnekleri yineler. Ayrıca o kesimden biraz daha kısadır.
\end{editorial}

Sonlu bir kümeyi yalnızca !!{element}larını numaralandırarak belirtebiliriz.
Şöyle bir küme tanımladığımızda bunu yaparız:
\[
  A = \{a_1, a_2, \ldots, a_n\}.
\]
$a_1$, \dots, $a_n$ !!{element}larının birbirinden farklı olduğunu
varsayarsak bu, $A$ ile ilk $n$ doğal sayı $0$, \dots, $n-1$ arasında bir
!!{bijection} verir. Tersine, her sonlu kümenin yalnızca sonlu sayıda
!!{element}ı olduğundan her sonlu küme böyle bir eşleme içine sokulabilir.
Başka bir deyişle $A$ sonluysa $A$ ile $\{0,\dots,n-1\}$ arasında bir
!!{bijection} vardır; burada $n$, $A$'nın !!{element}larının sayısıdır.

Belirli türde sonsuz kümelere izin verirsek bazı sonsuz kümelerin de
numaralandırılmasına izin veririz. Sonsuz bir kümenin, kendisiyle $\Nat$'ın
tamamı arasındaki bir !!{bijection} ile numaralandırıldığını söyleyerek bunu
kesinleştirebiliriz.

\begin{defn}[Numaralandırma, küme kuramsal]
Bir $A$ kümesinin \emph{numaralandırması}, görüntü kümesi $A$, tanım kümesi
  ise doğal sayıların bir başlangıç parçası $\{0,1,\ldots,n\}$ ya da doğal
sayılar kümesinin tamamı $\Nat$ olan bir !!{bijection}dur.
\end{defn}

\begin{explain}
\emph{Numaralandırma} sözcüğünün bu kullanımının sezgisel bir temeli vardır.
Bir $A$ kümesini numaralandırdığımızı söylemek, $A$ kümesinin elemanlarını
saymamızı sağlayan bir !!{bijection} $f$ bulunduğunu söylemektir. $0$'ıncı
eleman $f(0)$, $1$'inci $f(1)$, \ldots, $n$'inci $f(n)$,\ldots{} olur.
\footnote{Evet, $0$'dan sayıyoruz. Elbette~$1$'den de başlayabilirdik. Bunun
önemli bir farkı olmazdı; yalnızca~$\Nat$ yerine~$\PosInt$ kullanmamız
gerekirdi.} Bunun gerekçesi şu eklemeyle daha da açık hâle getirilebilir:
\end{explain}

\begin{defn}
  \ollabel{defn:enumerable}
  Bir $A$ kümesi, ancak ve ancak $A=\emptyset$ ise ya da $A$'nın bir
  numaralandırması varsa !!{enumerable}dir. $A$ !!{enumerable} değilse
  !!{nonenumerable} olduğunu söyleriz.
\end{defn}

\begin{explain}
Öyleyse bir küme, ancak ve ancak boşsa ya da !!{element}ları bir
numaralandırmayla sayılabiliyorsa !!{enumerable}dir.
\end{explain}

\begin{ex}
Doğal sayıların bir numaralandırması, $\Id{\Nat}(n)=n$ ile verilen özdeşlik
fonksiyonu $\Id{\Nat}\colon\Nat\to\Nat$'tır. \emph{Pozitif} doğal sayılar
kümesi $\Nat^+=\Nat\setminus\{0\}$'ın bir numaralandırması ise $g(n)=n+1$,
yani ardıl fonksiyonudur.
\end{ex}

\begin{prob}
Bir $A$ kümesinin !!{enumerable} olmasının, ya $A=\emptyset$ olmasına ya da
bir !!{surjection} $f\colon\Nat\to A$ bulunmasına eşdeğer olduğunu gösterin.
$A$'nın !!{enumerable} olmasının bir !!{injection} $g\colon A\to\Nat$
bulunmasına eşdeğer olduğunu gösterin.
\end{prob}

\begin{ex}
\begin{align*}
  f(n) & = 2n \text{ ve}\\
  g(n) & = 2n+1
\end{align*}
ile verilen $f\colon\Nat\to\Nat$ ve $g\colon\Nat\to\Nat$ fonksiyonları
sırasıyla çift doğal sayıları ve tek doğal sayıları numaralandırır. Fakat
hiçbiri !!{surjective} değildir; dolayısıyla hiçbiri $\Nat$'ın bir
numaralandırması değildir.
\end{ex}

\begin{prob}
$1$, $4$, $9$, $16$, \dots{} kare sayılarının bir numaralandırmasını
tanımlayın.
\end{prob}

\begin{ex}
$\lceil x\rceil$, $x$'i en yakın üst tam sayıya yuvarlayan \emph{tavan}
fonksiyonu olsun. O zaman
\[
  f(n) = (-1)^{n} \left\lceil\tfrac{n}{2}\right\rceil
\]
ile verilen $f\colon\Nat\to\Int$ fonksiyonu tam sayılar kümesi $\Int$'i
şöyle numaralandırır:
\[
\begin{array}{c c c c c c c c}
f(0) & f(1) & f(2) & f(3) & f(4) & f(5) & f(6) & \dots \\ \\
\big\lceil \tfrac{0}{2} \big\rceil & -\big\lceil \tfrac{1}{2}\big\rceil &  \big\lceil \tfrac{2}{2} \big\rceil & -\big\lceil \tfrac{3}{2} \big\rceil & \big\lceil \tfrac{4}{2} \big\rceil  & -\big\lceil \tfrac{5}{2}\big\rceil & \big\lceil \tfrac{6}{2} \big\rceil & \dots \\ \\
0 & -1 & 1 & -2 & 2 & -3 & 3& \dots
\end{array}
\]
$f$'nin pozitif ve negatif tam sayılar arasında ileri geri ``sıçrayarak''
$\Int$ değerlerini nasıl ürettiğine dikkat edin. $f$'yi durumlara göre şöyle
tanımlanmış olarak da düşünebilirsiniz:
\[
f(n) = \begin{cases}
  \frac{n}{2} & \text{$n$ çift ise}\\
  -\frac{n+1}{2} & \text{$n$ tek ise.}
  \end{cases}
\]
\end{ex}

\begin{prob}
$A$ ve $B$ !!{enumerable} ise $A\cup B$'nin de !!{enumerable} olduğunu
gösterin.
\end{prob}

\begin{prob}
$A_1$, $A_2$, \dots, $A_n$ kümelerinin hepsi !!{enumerable} ise
$A_1\cup\dots\cup A_n$'nin de !!{enumerable} olduğunu $n$ üzerinde
tümevarımla gösterin.
\end{prob}

\end{document}
